\section{Introduction}\label{intro}
In epidemiology, we are often interested in quantifying duration distributions,
\eg durations of transient health states or exposure to risk/protective factors.
Such distributions can help
characterize populations' cumulative exposure \cite{Ben-Shlomo2023}, and
inform inputs to various kinds of simulation models \cite{Garnett2011}.
\par
In HIV epidemiology, sex work
--- selling or buying sex for money or equivalent ---
typically reflects a transient period of
elevated risk for HIV acquisition and transmission due to
intersecting structural, behavioural, and biological factors
\cite{Scorgie2012,Shannon2015,Baral2015msw,Carael2006,Sanders2013cli}.
Unmet HIV prevention and treatment needs during this period
have long been recognized as drivers of epidemic emergence and persistence,
especially through insights from mathematical models of HIV transmission
\cite{Boily2002,Watts2010,Silhol2026}.
While entry into sex work is driven primarily by economic and social instability,
exit is further mediated by aging, stable partnering, stigma, and violence
\cite{Ingabire2012,Bowen2011,Lora2025}.
% JK: is it worth adding analogous drivers for clients? e.g. Sanders2013cli
%     It felt awkward when I tried.
\par
Durations selling or buying sex are thus
key determinants of HIV epidemiology, and key inputs to HIV transmission models,
particularly those aiming to compare the impacts of different prioritized services.
For example, 10\% HIV prevalence among sex workers
in a context with shorter typical durations / rapid turnover (\eg 1 year)
suggests a high HIV incidence rate, whereas 10\% HIV prevalence
in a context with longer typical durations / slower turnover (\eg 10 years)
suggests a relatively lower HIV incidence rate,
since prevalence has more time to accumulate \cite{Knight2020}.%
\footnote{Assuming HIV prevalence among new sex workers is negligible.}
% JK: I know publishers hate footnotes,
%     but it's only because they refuse to use LaTeX :)
%     https://history.stackexchange.com/q/46506/
Turnover can also moderate levels of
program / service coverage among sex workers (\eg PrEP),
where rapid turnover could result in
a constant influx of new sex workers to reach, and
potentially high loss-to-follow-up among those exiting sex work,%
\footnote{Questions remain regarding to what degree
  those exiting sex work remain engaged with programs.}
whereas slower turnover could make it easier to attain higher levels of coverage.
Finally, duration in sex work can further influence onward transmission risk,
where longer durations theoretically allow
more onward transmission following acquisition (in the absence of treatment).
% JK: I considered adding something to the above
%     about sex work becoming a net "sink" vs "source" of infections
%     i.e. net # infections going in vs coming out of sex work overall.
%     but it felt too awkward & tangential.
\par
Ideally, context-specific estimates of sex work duration could be obtained
from representative longitudinal cohorts with minimal loss-to-follow-up, or
from representative samples of "retired" sex workers.%
\footnote{Throughout:
  we often refer to "sex workers" for brevity,
  but our statements usually apply similarly to their clients who buy sex;
  we also intentionally use gender-agnostic terminology,
  since men, women, and transgender individuals can all sell and buy sex.}
Unfortunately, such studies are rare,
and most estimates derive from cross-sectional surveys of
those actively selling or buying sex \cite{Fazito2012}.
In their global scoping review of HIV risk behaviour durations,
\citet{Fazito2012} considered three methods to estimate
mean duration selling/buying sex:
1. mean total duration among those finished selling/buying sex;
2. reciprocal of the proportion who started selling/buying sex within the past year; and
3. double the mean/median reported duration among those actively selling/buying sex.
Among the 118 estimates of duration selling or buying sex in \cite{Fazito2012}:
the vast majority used method~3 (86\%),
a small faction used method~2 (14\%),
and none used method~1.
Yet, methods 2 and 3 have their own limitations (Table~\ref{tab:meth}),
especially overlooking so-called "length-biased sampling" \cite{Kafadar2009},
as we will discuss further below.
Moreover, studies also commonly report the proportions of respondents
selling/buying sex for pre-specified duration ranges, \eg \{0--2, 2--5, 5+\} years,
but it's not clear if/how such data could be used with methods 2 or 3.
As one recent HIV modelling study noted
when specifying plausible duration parameter bounds:
\emph{"We multiply the lower bound by 0.5 and upper bound by 2
  because of uncertainty in how overall duration of sex work
  relates to current duration of sex work"} \cite{Stone2021}.
\begin{table}
  \caption{Methods of estimating mean duration selling sex
    suggested and reviewed in \cite{Fazito2012}}
  \label{tab:meth}
  \input{tab.meth}
\end{table}
% ------------------------------------------------------------------------------
\paragraph{Objectives}
So, we sought to offer updated guidance on estimating duration in sex work
from different kinds of commonly available data.
Specifically, our objectives were to:
\begin{enumerate}
  \item \label{obj:1}
        Describe the exact analytic relationship between
        the distribution of reported durations selling sex in cross-sectional studies, and
        the distribution of total durations selling sex in the population overall.
  \item \label{obj:2}
        Apply a Bayesian model of the above relationships to infer
        mean total durations selling sex in the population overall
        from commonly available aggregate data,
        and compare resulting estimates with previously published estimates.
\end{enumerate}
While we focus on sex work as an illustrative example,
our insights likely extend to other transient health/risk states
relevant to compartmental modelling or epidemiological study.
The remainder of our paper is structured as follows:
In \sref{pre}, we clarify some key concepts and assumptions.
In \sref{math}, we derive the analytic relationship noted above.
In \sref{app}, we introduce and apply our Bayesian model,
  comparing mean total durations with previous estimates.
In \sref{disc}, we conclude with discussion of implications and future work.

